\documentclass[11pt,letterpaper]{article}
\usepackage[T1]{fontenc}
\usepackage[utf8]{inputenc}
\usepackage{lmodern}
\usepackage[margin=0.88in,headheight=15pt]{geometry}
\usepackage{amsmath,amssymb,amsthm,booktabs,longtable,array,enumitem,xcolor,microtype,fancyhdr,hyperref}
\definecolor{navy}{HTML}{17354A}
\definecolor{teal}{HTML}{176B73}
\definecolor{graytext}{HTML}{4F5961}
\hypersetup{colorlinks=true,linkcolor=navy,urlcolor=teal,citecolor=teal,
pdftitle={Exact table optimization: model, algorithm, optimality and limitations},
pdfauthor={}}
\pagestyle{fancy}\fancyhf{}
\fancyhead[L]{\small\color{graytext}Exact table optimization}
\fancyhead[R]{\small\color{graytext}Model and implementation limits}
\fancyfoot[L]{\small\color{graytext}16 October 2026}
\fancyfoot[R]{\small\thepage}
\renewcommand{\headrulewidth}{0.3pt}
\setlength{\parindent}{0pt}\setlength{\parskip}{5pt}
\setlength{\emergencystretch}{2em}
\setlist{itemsep=3pt,topsep=5pt}
\newcommand{\N}{\mathbb N}\newcommand{\Z}{\mathbb Z}\newcommand{\Q}{\mathbb Q}
\newcommand{\PP}{\mathsf{PhaseProduct}}
\newcommand{\OPT}{\operatorname{OPT}}
\newcommand{\bits}{\operatorname{bitlength}}
\newcommand{\val}{\operatorname{val}_2}
\newcommand{\cost}{\mathcal T}
\newcommand{\Plain}[1]{\par\textbf{\color{navy}Plain language.} #1\par}
\newcommand{\STE}[1]{\par\textbf{\color{teal}ASD-STE100 controlled English.} #1\par}
\newcommand{\Formal}[1]{\par\textbf{Formal statement.} #1\par}
\newtheorem{theorem}{Theorem}[section]
\newtheorem{lemma}[theorem]{Lemma}
\newenvironment{scopebox}{\begin{quote}\small\color{navy}\hrule\vspace{6pt}}{\vspace{6pt}\hrule\end{quote}}
\begin{document}
\hypersetup{pageanchor=false}
\begin{titlepage}
\vspace*{0.45in}
{\large\color{teal}IMPLEMENTATION REFERENCE\par}
\vspace{0.25in}
{\Huge\bfseries\color{navy}Exact table optimization\par}
\vspace{0.15in}
{\LARGE Model, algorithm, optimality,\\and limitations\par}
\vspace{0.35in}
{\large Kahanamoku-Meyer and Yao: Appendix F.2 points\\
Existing common-child-width PhaseProduct planner\par}
\vspace{0.25in}16 October 2026
\vspace{0.3in}
\begin{scopebox}
\textbf{Purpose.} Document the existing exhaustive exact algorithm for implementation and practicality tests. This document proposes no search improvement.

\textbf{Guarantee.} A completed, correct execution returns the minimum recursive T-count in the precisely defined model, a recursion sequence, and one realizing table per used arity.

\textbf{Limit.} The established time bound is
\[
2^{O(K^4\log K)}\operatorname{poly}(n).
\]
This is a theoretical upper bound, not a completion-time estimate.
\end{scopebox}
Each of the twelve limitations appears in three forms: ordinary language,
controlled technical English, and a formal mathematical statement.
The model, search procedure, optimality argument and complexity
accounting precede the itemization.

\textbf{Provenance.} KMY notation and point selection come from
arXiv:2403.18006v4 \cite{kmy}. The scanner, asymmetric charges, rounded
T-count objective, shared-table restriction and exhaustive optimization
theorems belong to the local research model \cite{threshold,outer,profile}.
They are not attributed to KMY.

\textbf{Verification status.} The full optimizer is not
machine formalized, and no large-arity execution is reported here.

\vfill
{\small\color{graytext}Standalone reference. Inputs, acceptance conditions, costs and stopping claims are defined here. References identify provenance and the full supporting proofs.}
\end{titlepage}
\pagenumbering{roman}
\hypersetup{pageanchor=true}
\begingroup
\small
\setlength{\parskip}{0pt}
\tableofcontents
\endgroup
\clearpage
\pagenumbering{arabic}

\section{Notation and exact model}\label{sec:model}
\subsection{Relation to KMY notation}
KMY use \(n\) for operand bitwidth, \(k\) for the number of chunks,
\(q=2k-1\) for PhaseProduct evaluations, and \(w_\ell\) for evaluation
points. Their operation is \cite[Eq.~(10)]{kmy}
\[
\PP(\phi)\lvert x\rangle\lvert z\rangle
=e^{i\phi xz}\lvert x\rangle\lvert z\rangle.
\]
Their coefficient vectors list high-degree coefficients first:
\(\mathbf x=(x_{k-1},\ldots,x_0)^{\mathsf T}\), and similarly
\(\mathbf z\). Their integer matrices \(A'\) and \(A''\) give
the Appendix B form
\[
\phi xz=\sum_{\ell=0}^{q-1}\phi_\ell
(A'\mathbf x)_\ell(A'\mathbf z)_\ell,\qquad
\boldsymbol\phi^{\mathsf T}
=\phi\,\mathbf e_{b_{\rm rad}}^{\mathsf T}(A'')^{-1}.
\tag{KMY--PP}
\]
Here \(\mathbf e_{b_{\rm rad}}\) has length \(q\). At current width \(m\),
our unequal-chunk convention uses \(b_{\rm rad}=2^{\lfloor m/k\rfloor}\).
The identity specifies the required linear forms; the table optimizer
does not search over the phase coefficients.

\begin{longtable}{@{}p{0.19\textwidth}p{0.75\textwidth}@{}}
\toprule Symbol & Meaning\\ \midrule
\(n\ge1,\ K\ge2\)&Root operand bitwidth and maximum allowed arity.\\
\(m\)&Current operand bitwidth; never an evaluation point.\\
\(k,\ q=2k-1\)&Arity and number of required PhaseProduct markers.\\
\(w_\ell,A',A''\)&KMY evaluation point and integer matrices.\\
\(P,X\)&Complete instruction table and current integer row matrix.\\
\(\gamma(P)\)&Positive additions plus three times negative additions plus negations.\\
\((\alpha,\beta)\)&Canonical two-parameter historical-width profile.\\
\(W_P(m)\)&Common child width specified by the scanner.\\
\(C,\delta,\Lambda\)&Charge budget, profile difference and width cutoff.\\
\(N_i,R\)&Call multiplicity and terminal rotation count \(N_Dm_D^2\).\\
\(\cost,U\)&Recursive T-count and the cost of stopping at the root.\\
\bottomrule
\end{longtable}

The program indexes coefficients in the opposite order. Let \(J\)
reverse coordinates and put
\(\mathbf x_\uparrow=J\mathbf x=(x_0,\ldots,x_{k-1})^{\mathsf T}\).
Program row \(X_{i,*}\) denotes \(X_{i,*}\mathbf x_\uparrow\).
Its target is the corresponding row of \(A'J\).
This reversal changes notation, not physical operations.

\subsection{Fixed points and literal target rows}
The points are the first \(q\) entries of the F.2 sequence
\[
0,\ \infty,\ -1,\ 1,\ -\tfrac12,\ \tfrac12,\ -2,\ 2,\
-\tfrac14,\ \tfrac14,\ -4,\ 4,\ldots .
\tag{1}
\]
For successive \(\omega=2,4,8,\ldots\), the next block is
\((-\omega^{-1},\omega^{-1},-\omega,\omega)\).
This selects a set, not a mandatory marker order.
KMY Appendix G examples can use different point sets
\cite[App.~F.2 and G]{kmy}.

For coefficient index \(j=0,\ldots,k-1\), define
\[
\begin{aligned}\relax
[v_k(a)]_j&=a^j &&(a\text{ an integer}),\\
[v_k(1/c)]_j&=c^{k-1-j}
 &&(c\in\{\pm2,\pm4,\ldots\}),\\
v_k(\infty)&=e_{k-1}.
\end{aligned}\tag{2}
\]
Use \(0^0=1\). A negative reciprocal keeps its negative denominator
\(c\). Thus \(v_4(-1/2)=(-8,4,-2,1)\), not \((8,-4,2,-1)\).
Sign equivalence is not literal equality. Every target is primitive:
its entries have no common factor 2.

For \(k=3\), the five required rows are
\[
\begin{array}{c|ccccc}
w_\ell&0&\infty&-1&1&-1/2\\ \hline
v_3(w_\ell)&(1,0,0)&(0,0,1)&(1,-1,1)&(1,1,1)&(4,-2,1).
\end{array}
\]
Arbitrary evaluation points are not inputs to this problem.

For full matrix definitions, let \(J_t\) reverse \(t\) coordinates.
The rows in descending order are
\[
 A'_{\ell,*}=v_k(w_\ell)J_k,\qquad
 A''_{\ell,*}=v_q(w_\ell)J_q,\qquad
 \mathbf e_{b_{\rm rad}}=(b_{\rm rad}^{q-1},\ldots,b_{\rm rad},1)^{\mathsf T}.
\]
The rule (2) defines \(v_q\) by replacing its vector length \(k\) with
\(q\), while keeping the same selected points \(w_\ell\).
For infinity this selects the leading coefficient.
The distinct points make the scaled interpolation matrix \(A''\)
invertible. No numerical interpolation approximation is needed to
define the model.


\subsection{Legal instruction tables}\label{sec:dsl}
Start with \(X=I_k\). All \(k\) rows are writable.
For distinct \(d,h\), sign \(\varepsilon=\pm1\), and integer \(s\ge0\):
\[
\begin{array}{lll}
\mathrm{add}(d,h,\varepsilon,s):
 &X_d\leftarrow X_d+\varepsilon2^sX_h,
 &\text{charge }1\text{ or }3;\\
\mathrm{left}(d,s):
 &X_d\leftarrow2^sX_d,&\text{charge }0;\\
\mathrm{right}(d,s):
 &X_d\leftarrow X_d/2^s,&\text{charge }0;\\
\mathrm{neg}(d):
 &X_d\leftarrow-X_d,&\text{charge }1;\\
\mathrm{mark}(d,\ell):
 &\text{require }X_d=v_k(w_\ell),&\text{charge }0.
\end{array}\tag{3}
\]
An addition costs 1 when \(\varepsilon=1\) and 3 when
\(\varepsilon=-1\). A right shift is legal only when every entry of
the selected row is divisible by \(2^s\). A marker changes no row.

Acceptance requires every target exactly once and final matrix
exactly \(I_k\), with the original row indices and signs.
Every intermediate row must be integral. Operations are invertible
over \(\Q\), so no reachable row is zero.

Arbitrary legal helper rows, helper sharing and marker permutations
are allowed. There is no separate reset, helper-shape, determinant-size
or instruction-count restriction. A finite minimum-witness bound is
a theorem about this universe, not an extra constraint on it.

\subsection{Historical width scanner}\label{sec:scanner}
At operand width \(m=kL+r\), \(0\le r<k\), initial chunk widths are
\((L,\ldots,L,L+r)\). The scanner updates are
\[
\begin{array}{ll}
\mathrm{add}(d,h,\varepsilon,s):&u_d\leftarrow1+\max(u_d,u_h+s),\\
\mathrm{left}(d,s):&u_d\leftarrow u_d+s,\\
\mathrm{right}(d,s):&u_d\leftarrow\max(0,u_d-s),\\
\mathrm{neg}(d):&u_d\leftarrow u_d+1.
\end{array}\tag{4}
\]
Markers change nothing. Let \(H(P;u_{\rm in})\) be the maximum of all
coordinates at all boundaries, including the initial boundary. Define
\[
W_P(m)=1+H(P;(L,\ldots,L,L+r)).\tag{5}
\]
This conservative specified scanner is not an exact signed-value
range analysis or a branch-specific KMY width calculation.

Put \(R_k=k-1\). Two complete scans determine the width function:
\[
\begin{split}
H_0&=H(P;(0,\ldots,0)),\\
H_R&=H(P;(0,\ldots,0,R_k)),\\
\alpha(P)&=H_0,\qquad \beta(P)=H_R-R_k.
\end{split}\tag{6}
\]
Every complete legal table obeys \cite{profile}
\[
0\le\beta\le\alpha\le\beta+k-1,\qquad
W_P(m)=\left\lfloor m/k\right\rfloor+1+
       \max(\alpha,(m\bmod k)+\beta).
\tag{7}\label{eq:width}
\]
Thus \((k,\gamma,\alpha,\beta)\) is an exact planner summary for a
completed table. It is not a sufficient search state for an
unfinished table.

For a canonical cap \((a,b)\), the raw historical tests are
\(H_0\le a\) and \(H_R\le k-1+b\). The second raw cap is not \(b\).
All peaks must include the full history.

\subsection{Exact recursive objective}\label{sec:objective}
A plan is a finite arity sequence \(\pi=(k_0,\ldots,k_{D-1})\),
\(2\le k_i\le K\), with one complete legal table \(P_k\) per used arity.
Set
\[
m_0=n,\quad N_0=1,\quad
m_{i+1}=W_{P_{k_i}}(m_i)<m_i,\quad
N_{i+1}=N_i(2k_i-1).
\tag{8}
\]
All calls at a level use this same arity, table and child width.
Repeated arities reuse the same table. The empty sequence is allowed.
No monotone-arity requirement is imposed.

For positive integer \(R\), define
\[
\begin{split}
\tau(R)
 &=\left\lfloor 3.075\log_2(100R)+0.86+\tfrac12\right\rfloor\\
 &=\left\lfloor
 \frac{\bits((100R)^{615})-1+272}{200}
 \right\rfloor .
\end{split}\tag{9}\label{eq:tau}
\]
The constants are exact rationals and rounding is
\(\lfloor x+1/2\rfloor\), not ties-to-even. The second line permits
exact evaluation without floating-point logarithms.

The repository's chosen PP benchmark is
\[
\boxed{\quad
\cost(\pi,\{P_k\})=
8\sum_{i=0}^{D-1}N_i(m_{i+1}-1)\gamma(P_{k_i})
+R(4+\tau(R)),\qquad R=N_Dm_D^2 .
\quad}\tag{10}
\]
Define \(\OPT(n,K)\) as the minimum of (10) over all feasible plans
and legal tables. Formula (10) is not KMY's original
\(CR_\phi\)-count estimator. Exact optimization of this rounded cost
does not mean exact Clifford+T synthesis of arbitrary phase angles.

\begin{scopebox}
\textbf{Output.} The algorithm returns the minimum cost, a realizing
arity sequence, and one explicit marked instruction table per used
arity. Replay must establish all widths, multiplicities, charges and
the terminal cost. A feasible table alone does not certify a minimum.
\end{scopebox}

\section{How the exhaustive algorithm finds tables}\label{sec:algorithm}
\subsection{Minimum charge under a width cap}
For a canonical pair \((a,b)\), let
\[
c_{\min}(k,a,b)=
\min\{\gamma(P):P\text{ complete and legal},\
 \alpha(P)\le a,\ \beta(P)\le b\},
\tag{11}
\]
with \(+\infty\) for an empty set. For
\(\delta\in\{0,\ldots,k-1\}\), define
\[
\theta(k,\delta,C)=
\min\{a\ge\delta:c_{\min}(k,a,a-\delta)\le C\}.
\tag{12}
\]
Set \(\min\varnothing=+\infty\), so an infeasible ray has infinite
threshold.
A used table has \(\alpha\le n-2\): its child width is at least
\(\alpha+1\) and at most \(n-1\). Therefore only thresholds at most
\(A_n=n-2\) matter. At \(n=1\), root stop is the only plan.

The oracle returns the exact threshold and a witness when
\(\theta\le A_n\), or reports \(\theta>A_n\).
The latter is not a global infinity claim unless the full sufficient
cutoff was searched.

\subsection{Why the search is finite}\label{sec:bounds}
The reviewed algorithm uses three established facts
\cite{threshold,normalization,outer}. The following statements and
proof summaries identify the mathematical dependencies.

\begin{lemma}[Charge ceiling and incumbent]\label{lem:charge}
Every feasible cap has a minimum-charge witness of charge at most
\(27k^3\). An explicit paired F.2 constructor has charge
\(c_0\le6k^2\) and profile \((a_0,b_0)\) with \(a_0,b_0\le8k^2\).
\end{lemma}
The paired constructor groups opposite points. It prepares even-
and odd-degree coefficient sums in two existing rows, emits the two
literal target sums using addition/subtraction, and reverses the
preparation. It forms and reverses an unpaired final target directly.
Negative-reciprocal signs are included. This construction supplies
an incumbent, not a restriction on candidate tables.

For the charge ceiling, the potential
\(\sum_j u_j-\log_2|\det X|\) increases by at least one at each
charged operation and is unchanged by exact shifts. Identity endpoints
therefore imply \(\gamma\le3ka\) under zero-seed cap \(a\).
For fixed \(\delta\), put
\(a_*=\max(a_0,b_0+\delta)<9k^2\).
Below \(a_*\), the potential bounds the minimum by \(27k^3\);
at and above \(a_*\), the constructor costs at most \(6k^2\).
The constructor bounds are proved in the retained source package,
not inferred from its finite tests.

\begin{lemma}[Primitive normalization]\label{lem:primitive}
Every complete legal table has a replacement preserving signed
operation counts and literal marker order, with no greater historical
width in either scan. Rows are primitive between charged macros.
A minimum witness needs \(O(k^3)\) primitive instructions.
\end{lemma}
Write \(X_i=2^{g_i}Y_i\), with \(Y_i\) primitive, and replace width
\(u_i\) by \(u_i-g_i\). The row-mass invariant
\(\|X_i\|_1\le2^{u_i}\) makes this nonnegative. Absorb isolated shifts
into the next addition. A negative relative source exponent becomes
a destination left shift. After addition, divide by the row's common
power of two. The invariant proves that every transformed transient
peak is no greater than the corresponding old peak.
Primitive literal targets and identity endpoints are preserved.
A conservative length bound is \(3A+\nu+q\), for \(A\) additions
and \(\nu\) negations, plus four optional padding instructions.

\begin{lemma}[Threshold saturation]\label{lem:saturation}
For \(C<c_0\), a finite threshold satisfies
\[
\theta(k,\delta,C)\le S(k),\qquad
S(k)=(6k^2+1)(6k^2)!(8k^2)^{6k^2}
=2^{O(k^2\log k)}.
\tag{13}
\]
For \(C\ge c_0\), it satisfies \(\theta(k,\delta,C)\le a_*<9k^2\).
\end{lemma}
In the low-budget proof, each scalar coefficient of a normalized
operation pattern is a signed sum of powers of two. Finite alternatives
for exponent order, small gaps and separated cancellation clusters
give integer linear constraints. They enforce integrality at every
prefix, all literal markers, identity restoration and both histories.
At a primitive boundary, if either relative pre-addition exponent
is positive, the other exponent is zero and parity makes the sum
primitive. Otherwise both pre-addition exponents are zero and only
the normalization-division exponent can vary.
Thus a fixed addition template has at most one exponent variable;
the cap adds one more.
There are at most \(6k^2\) variables in this budget range, with
coefficient and constant bounds \(8k^2\).

The small-integer-point lemma used by the proof is
\[
\begin{gathered}
 B\in\Z^{s\times t},\quad d\in\Z^s,\quad t,M\in\Z_{\ge1},\\
\{x\ge0:Bx\le d\}\cap\Z^t\ne\varnothing,\quad
 |B_{ij}|,|d_i|\le M\\
\Longrightarrow\quad
\text{some feasible integer }y\text{ has }
0\le y_i\le(t+1)t!M^t.
\end{gathered}\tag{14}
\]
Determinant bounds on vertices and integer recession rays, followed
by removal of integer ray multiples, prove (14), including
lower-dimensional systems via homogenization.
This gives (13). The high-budget bound uses the explicit constructor.
The symbolic argument proves a cutoff; the executable search below
does not need to enumerate the symbolic systems.

\subsection{Normalized macro enumeration}
For each \((k,\delta,C)\), set
\[
\Lambda=
\begin{cases}
\min(n-2,S(k)),&C<c_0,\\
\min(n-2,\max(a_0,b_0+\delta)),&C\ge c_0.
\end{cases}\tag{15}
\]
If \(n-2<\delta\), report \(\theta>n-2\).
Budget zero is infeasible for \(k\ge2\), because a nonbasis target
must be created.

A state contains
\[
(Y,u^0,u^R,H_0,H_R,S,c),
\tag{16}
\]
where \(Y\) is the primitive integer matrix, \(S\) is the marked-target
set and \(c\) is weighted charge. Store predecessor information for
reconstruction. Initially \(Y=I_k\), \(c=0\), and the width seeds are
those of (6). Initial targets can be marked immediately.

Enumerate these choices:
\begin{enumerate}
\item \textbf{Addition.} Choose \(d\ne h\), \(\varepsilon=\pm1\)
and \(p_d,p_h\ge0\), with
\(\min(p_d,p_h)=0\), \(p_d,p_h\le\Lambda-1\). Form
\[
Z=2^{p_d}Y_d+\varepsilon2^{p_h}Y_h,\qquad
t=\min_{j:Z_j\ne0}\val(|Z_j|).
\]
Before division, test
\[
v^s=1+\max(u_d^s+p_d,u_h^s+p_h),\qquad s\in\{0,R\},
\]
against raw caps \(\Lambda\) and \(k-1+\Lambda-\delta\).
Record these peaks. If valid, set
\(Y_d=Z/2^t\), \(u_d^s=v^s-t\), and add charge 1 or 3.
Expand the macro as a destination left shift when needed, the
permitted addition, and an exact right shift when \(t>0\).

\item \textbf{Negation.} Negate a selected row, increase its scanner
width by one in both scans, check both caps, and add charge one.

\item \textbf{Mark and accept.} At each primitive boundary, emit
every newly available required literal marker once.
Choose a deterministic row if needed.
Markers are free and value-neutral, and their order is unrestricted.
Accept only states with all targets marked and \(Y=I_k\).
\end{enumerate}

Reject any prefix of weighted charge greater than \(C\).
Each macro has positive charge, so depth is at most \(C\).
Depth is not a substitute for the weighted-budget test.

For an accepted word \(P\), its least cap on this ray is
\[
a_{\rm ray}(P)=\max(\delta,\alpha(P),\beta(P)+\delta).
\tag{17}
\]
Return the minimum of these values with its explicit word.
If no candidate exists, return \(\theta>n-2\).
Only exhaustion of the full untruncated sufficient cutoff establishes
that no finite threshold exists.

\subsection{Construct the minimum-charge catalogue}
There are \(k\) differences and \(27k^3+1\) charge budgets.
Thresholds determine
\[
c_{\min}(k,a,a-\delta)
=\min\{C:\theta(k,\delta,C)\le a\}.
\tag{18}
\]
A witness below cap \((a,b)\) can be given exactly that canonical
profile by prepending
\[
\mathrm{left}(0,a),\ \mathrm{right}(0,a),\
\mathrm{left}(k-1,b),\ \mathrm{right}(k-1,b).
\tag{19}
\]
This prefix has zero charge. It restores the matrix and current
widths, raises only historical peaks, and adds no markers.
Hence cap-minimum and exact-profile minimum charges agree.

The catalogue must retain the needed charge--width choices.
Keeping only the cheapest table per arity is insufficient.
An offset interval also cannot automatically be replaced by its
smallest endpoint: optimized continuation cost need not be monotone
in width.

\section{Why the answer is optimal}\label{sec:optimal}
\subsection{Local optimality}
\begin{theorem}[Threshold and cap-minimum realization]
The finite macro search returns the exact truncated threshold with
a legal witness. Equations (18)--(19) give a minimum-charge table
at each requested exact canonical profile.
\end{theorem}
\begin{proof}
Every returned word is sound by exact replay of primitives, markers,
caps and identity restoration. Conversely, a genuine witness has a
no-wider, same-charge normalized representative by
Lemma~\ref{lem:primitive}. Lemma~\ref{lem:saturation} supplies the
sufficient finite range whenever a threshold is finite.
Every macro of that representative occurs in the search alphabet,
and its total length is within the charge-depth bound.
Earliest marking loses no word because targets are primitive and
marker order is unrestricted. Therefore a minimizing representative
is considered. Taking the least ray cap gives the threshold.
Taking the least successful budget gives (18).
Padding gives the exact requested profile without additional charge.
\end{proof}
This is minimum charge \emph{at a specified profile}. A recursively
selected table need not have minimum charge without a width constraint.

\subsection{Global recursive optimization}
Let root-stop cost be
\[
U=n^2(4+\tau(n^2)).
\tag{20}
\]
An optimum costs at most \(U\). Since positive-width terminal cost
is at least its multiplicity, every competitive plan satisfies
\[
Q=N_D=\prod_i(2k_i-1)\le U,\qquad
D\le\lfloor\log_3 U\rfloor.
\tag{21}
\]
The retained exact outer procedure \cite{outer} does the following:
\begin{enumerate}
\item Enumerate arity sequences with \(Q\le U\), decorated by residues
\(r_i=m_i\bmod k_i\), profile differences \(\delta_{k_i}\), and charges
\(c_{k_i}\). Repeated arities must agree on difference and charge.
\item Along a fixed \((k,\delta)\), the minimum charge is nonincreasing
in the offset \(a\). A fixed charge label occupies one integer interval.
Assign one integer offset variable \(a_k\) in that interval to each
distinct used arity.
\item For each terminal width \(z=1,\ldots,n\), enforce exact recursion,
strict contraction and the objective by exact integer feasibility.
Find the least budget and recover the offsets.
\item Compare every feasible result with root stop.
Reconstruct one padded witness per used arity and reuse it everywhere.
\end{enumerate}

The width elimination in step 3 is explicit. Set
\(P_0=1\), \(P_{i+1}=P_i k_i\),
\(h_i=1+\max(0,r_i-\delta_{k_i})\), and define affine expressions
\[
x_0=n,\qquad
x_{i+1}=x_i+P_{i+1}(a_{k_i}+h_i)-P_i r_i,\qquad
m_i=x_i/P_i.
\tag{22}
\]
Impose \(x_D=P_Dz\). All later increments are divisible by
\(P_{i+1}\), so reduction modulo \(P_{i+1}\) enforces the required
integer width and residue at every earlier level.
Strict contraction is
\[
k_i x_i-x_{i+1}\ge P_{i+1}.
\tag{23}
\]
For fixed \(z\), the leaf cost is constant and the remaining cost is
affine in the offsets. Multiplication by \(P_D\) clears denominators.
Integer budget queries recover the exact minimum.

\begin{theorem}[Global exactness in the declared model]
Complete local threshold results and exact outer optimization return
\(\OPT(n,K)\), a recursion sequence and a realizing shared-table family.
\end{theorem}
\begin{proof}
Take any feasible plan of cost at most \(U\). Every used offset
satisfies \(a\le n-2\). Replace each arity's table once by a minimum
catalogue witness with the \emph{same exact profile}, using padding.
Equation (7) preserves all widths, residues, multiplicities and
admissibility conditions, while charge does not increase.
Thus some catalogue plan is at least as good as any unrestricted
competitor. Bounds (21) include it in the outer search.
Shared offset variables enforce one table per arity.
Conversely, every accepted outer solution reconstructs legal tables
and an admissible plan. The two optimum values are equal.
\end{proof}

\clearpage
\subsection{Feasibility certificates and optimality certificates}
Replaying a table proves its feasibility and charge. Replaying a
plan proves its achievable T-count. Neither replay excludes a better
answer. An optimality claim also requires complete local searches
for the necessary smaller budgets/caps and complete outer optimization,
or sound certificates that cover every excluded alternative.

A timeout, a hash or a solver status string alone is not that proof.
No general polynomial-size independently checkable optimality
certificate is asserted. Search traces and infeasibility certificates
can themselves be large.

\subsection{Is verifying optimality efficient?}\label{sec:verification}
There are three different verification tasks.

\textbf{Verify a feasible answer.} Given explicit tables and a recursion
sequence, replay every operation using exact integers. Check target
coverage, identity restoration, both histories, repeated-arity
consistency, strict contraction and the score (10).
This takes polynomial time in the instruction count, arities and
integer bit lengths encountered. A successful replay proves an upper
bound on the optimum.

For a used table, its zero-seed history is at most \(n-2\).
The row-mass bound then gives \(O(n)\)-bit coefficients. The
minimum-charge ceiling and primitive normal form provide
\(O(k^3)\)-instruction witnesses per used arity, with shifts encoded
using \(O(\log(n+1))\) bits after padding. Together with
\(D=O(\log(n+1))\), this gives polynomial-size feasible witnesses in the measure
\(n+K+\bits(B)\). Thus the budget decision
\emph{does some feasible plan have cost at most \(B\)?} is in NP
when \(n,K\) are encoded in unary and \(B\) in binary.
If \(B\ge U\), the empty plan suffices; otherwise a feasible
plan has \(D\le\log_3 U\). This is not a binary-\(n\) certificate claim.

\textbf{Verify that no better answer exists.} For a feasible incumbent
of integer cost \(B\), optimality additionally requires
\[
 \nexists\,(\pi,\{P_k\})\text{ feasible with }
 \cost(\pi,\{P_k\})\le B-1.
\]
A cheaper plan can be certified by a feasible witness. The absence
of every cheaper plan is a different requirement. No uniform
polynomial-time verifier with a guaranteed polynomial-size
optimality certificate has been established here. The NP upper
bound for cheaper-plan existence does not provide such a certificate.

\textbf{Verify a supplied proof or exhaustive trace.} A sound checker
may run in time polynomial in the length of the supplied proof.
That proof or trace can itself be exponentially long in the problem
parameters. Polynomial verification in certificate length therefore
does not imply polynomial verification in the original input size.
A matching independently proved lower bound can make an individual
instance easy to certify, as in the \(k=2\) example.

For the current method, one general way to validate optimality is to
complete the exact searches and outer optimization described above.
Its available guarantee remains (26): polynomial in numeric \(n\)
for each fixed \(K\), but with no uniform polynomial guarantee when
\(K\) grows. The absence of a known efficient
optimality certificate is not a proof that efficient certification is
impossible, or that the optimization problem is outside P.

\clearpage
\subsection{A complete small example}
For \(k=2\), the targets are \(0,\infty,-1\). Consider
\[
\begin{array}{c|c|c|r}
\text{operation}&X_0&X_1&\text{charge}\\ \hline
\text{start; mark }0,\infty&(1,0)&(0,1)&0\\
X_0\leftarrow X_0-X_1;\ \text{mark }-1&(1,-1)&(0,1)&3\\
X_0\leftarrow X_0+X_1&(1,0)&(0,1)&4
\end{array}
\]
Raw peaks are \(2\) and \(3\), so the canonical profile is \((2,2)\).
At least two additions are needed to create an off-diagonal target
and restore identity. If any addition is negative, their total
charge is at least \(3+1=4\). If all additions are positive,
a negative coefficient requires a negation. Identity restoration
requires an even number of negations, since additions have determinant
one and shifts have positive determinant. There are then at least
two additions and two negations, again costing at least four.
Thus the displayed table is charge-optimal, both without a width cap
and under cap \((2,2)\).

Its child width is \(\lceil m/2\rceil+3\). At \(m=16\), the child
width is \(11\), and its local cost at multiplicity one is \(320\).
This does not establish that this step is part of the best
\emph{whole} recursive plan at \(n=16\).

\section{Computational complexity and storage}\label{sec:complexity}
The following are bit-operation bounds in numeric physical \(n\);
they are not polynomial-in-\(\log n\) claims.

\subsection{Local search}
The branching count is \(O(k^2(\Lambda+1))\), with depth at most \(C\).
Under the cap, coefficients have \(O(\Lambda)\) bits by the row-mass
bound, and exact replay costs polynomial time in \(k,\Lambda\).
Thus a query takes
\[
(k^2(\Lambda+1))^{O(C)}
\operatorname{poly}(k,\Lambda,\log(n+1)).
\tag{24}
\]
For \(C<c_0\), use \(C=O(k^2)\) and
\(\log\Lambda=O(k^2\log k)\).
For \(C\ge c_0\), use \(\Lambda=O(k^2)\) and \(C=O(k^3)\).
The respective parameter factors are
\(2^{O(k^4\log k)}\) and \(2^{O(k^3\log k)}\).
The uniform local bound is
\[
2^{O(k^4\log k)}\operatorname{poly}(\log(n+1)).
\tag{25}
\]
There are \(O(k^4)\) threshold queries per arity and \(O(K^5)\)
through arity \(K\).

\subsection{Outer search and full bound}
Each arity occurrence has at most \(28k^5\) decorations.
The retained counting argument gives fewer than \(3U^{10}\)
decorated sequences. If \(t\) distinct arities occur, then
\(t!\le Q\le U\), so \(t^{O(t)}=U^{O(1)}\).
The outer proof imports Kannan's exact integer-feasibility theorem
\cite[Theorem~5.4]{kannan}, in
\(t^{O(t)}\operatorname{poly}(\text{encoded input length})\) time.
Only the \(t\) offsets are integer variables; the widths in (22)
are expressions, not extra variables.

With terminal-width choices, budget searches, local queries and
reconstruction, the full bound is
\[
\boxed{2^{O(K^4\log K)}\operatorname{poly}(n).}\tag{26}
\]
This is fixed-parameter tractability in \(K\), under the
numeric-\(n\) convention. This bound does not establish uniform
polynomial time for growing \(K\). It supplies no lower bound that
excludes a faster algorithm.
In particular, the proved parameter factor is not
\(2^{O(K\log K)}\).

If a separate sufficient-cutoff theorem gave
\(K=O(\sqrt{\log n})\), substituting it into (26) would give
\[
2^{O((\log n)^2\log\log n)}\operatorname{poly}(n),
\]
a quasipolynomial upper bound, not a polynomial guarantee.
By contrast, a hypothetical factor \(2^{O(K\log K)}\) under that
same cutoff would be \(n^{o(1)}\), and would preserve a
\(\operatorname{poly}(n)\) total bound. Neither that smaller factor
nor a sufficient \(\sqrt{\log n}\) cutoff is assumed in this document.

\subsection{Time and memory interpretation}
The factor \(U^{10}\), hidden constants and local branching can be
large. Equation (26) gives no execution time at \(K=11\)--15 or
\(n=50{,}000\). The ceiling \(27k^3=91{,}125\) at \(k=15\)
is a sufficient search bound, not a prediction of optimal charge.

One matrix can require \(O(k^2(\Lambda+1))\) bits, in addition to
scanners, markers and reconstruction data. Cache size depends on
the implementation and can be much larger than one state.
The time theorem is not a practical RAM estimate; it does not require
all possible states to be stored at once.

The outer polynomial bound uses the stated integer-feasibility
dependence. A different solver can remain exact without inheriting
that bound. The local macro search itself needs no
integer-programming solver.

\clearpage
\section{Twelve limitations in three variants}\label{sec:limits}
Each item gives the same restriction in ordinary, controlled and
formal language. The controlled panels follow ASD-STE100 Issue 9
writing principles \cite{ste}, using the domain nouns defined here.
They are descriptive text. Their sentences have at most 20 words,
stricter than the descriptive 25-word limit. Complete conformity
also depends on dictionary meaning and part of speech.
No independent STE conformity certificate is claimed.
Equations and formal panels are outside this controlled layer.

\subsection{L1. Global optimality has a fixed search domain}
\Plain{The optimum covers all legal tables in Section~\ref{sec:model}.
The algorithm does not choose other evaluation points or search all
quantum arithmetic circuits. Arbitrary legal helpers and marker order
are allowed within the declared operation language.}
\STE{The algorithm uses one fixed set of evaluation points. The permitted row operations define the search domain. The global optimum applies only to this search domain. The algorithm does not select a different evaluation-point family.}
\Formal{Let \(\mathcal P_k\) be exactly the tables accepted by
(1)--(4). The minimum is over
\(\prod_{k\ {\rm used}}\mathcal P_k\) and admissible plans.
No equality is asserted with an optimum over other point sets,
extra registers, other primitives or arbitrary quantum circuits.}

\subsection{L2. The recursion model is restricted}
\Plain{All children share a width, and each arity has one table.
Every recursive step must reduce width. The guarantee does not
extend to independently optimized branches with different widths
or tables.}
\STE{Each recursion step uses one child width. The child width must be smaller than the parent width. Each arity uses the same table at all recursion steps. Other recursion models are outside the search domain.}
\Formal{Every feasible plan satisfies (8), with one \(P_k\) per arity.
It is a homogeneous-level recursion sequence. Branch-dependent
tables \(P_{k,i,\ell}\) and widths \(m_{i+1,\ell}\) are outside the
model. KMY's broader branch-dependent estimator is not identified
with this planner \cite[App.~F.1]{kmy}.}

\clearpage
\subsection{L3. The objective is one specified T-count}
\Plain{The optimizer minimizes the exact value of a rounded T-count
formula. It does not also minimize qubits, depth, physical runtime
or error-correction cost. Exact optimization of this formula is
not exact synthesis of arbitrary rotation angles.}
\STE{The objective is the recursive T-count. The objective does not include qubit count, circuit depth, or hardware time. A global optimum for T-count can have a higher cost for another objective.}
\Formal{The objective is precisely (10), using (9). In general,
\(\arg\min\cost\) is not contained in \(\arg\min G\) for another
resource \(G\). Charges \(1,3,1\), free shifts and the leaf cost
are fixed model data. The same recursive objective has not been
established here for PhaseTripleProduct.}

\subsection{L4. The arity cutoff limits the guarantee}
\Plain{The result is optimal among arities up to \(K\). A larger
arity can improve the answer. Practical observations of small
arities do not prove that a small cutoff preserves every optimum.
The search does not require decreasing arities.}
\STE{The input K gives the largest permitted arity. The global optimum can change when K increases. Tests at small arities do not give a proof that larger arities are unnecessary. The search permits arities in any order.}
\Formal{For \(K'\ge K\),
\(\OPT(n,K')\le\OPT(n,K)\).
Equality needs a sufficient-cutoff proof or exact comparison.
Neither \(K\le16\) nor \(k_{i+1}\le k_i\) is assumed.
The observed practical range \(11\)--15 is not an admissibility theorem.}

\clearpage
\subsection{L5. The time bound does not establish practicality}
\Plain{The exhaustive search can take extremely long. Hidden
constants and the growing arity factor prevent a reliable completion
estimate from the asymptotic bound. The polynomial is in operand
bitwidth, not the number of bits used to encode that width.}
\STE{The search can take a long time. The complexity bound does not give a completion time for a specific computer. The bound uses the operand bitwidth n. It does not use the bit length of n.}
\Formal{Equation (26) is an asymptotic upper bound in bit operations,
not a measured machine-time model. If \(b_n=\bits(n)\), then
\(\operatorname{poly}(n)\) can be \(2^{O(b_n)}\).
It does not establish uniform \(\operatorname{poly}(n,K)\) time,
and it does not establish non-P.}

\subsection{L6. The outer planner can remain costly}
\Plain{Table synthesis is only the local part. The planner must still
choose a recursion sequence, stopping level and consistent tables.
Its polynomial guarantee allows a large exponent and expensive
integer-feasibility calculations.}
\STE{The table search is only one part of the algorithm. The planner must also find a recursion sequence and consistent table choices. A polynomial bound can have a large exponent. The planner can therefore need much time.}
\Formal{The outer reduction permits fewer than \(3U^{10}\)
decorated sequences, \(n\) terminal widths and integer-budget
feasibility queries. Its dimension is \(t\), the number of distinct
arities, with \(t!\le U\). The imported \(t^{O(t)}\) algorithm is
a theoretical component, not a local benchmark. Another solver
requires separate runtime accounting even when it remains exact.}

\clearpage
\subsection{L7. Minimum local charge is not the recursive objective}
\Plain{A table with a higher charge can make its children narrower
and reduce total recursive cost. Keeping only the cheapest table
per arity can lose the global optimum. The required charge--width
alternatives must remain available, with one consistent table per arity.}
\STE{A table with a higher charge can give a smaller child width. The smaller width can reduce the recursive T-count. The planner must keep the required charge and width choices. The planner must use the same table when it uses the same arity again.}
\Formal{Minimizing \(\gamma(P_k)\) alone need not minimize (10).
Safe catalogue replacement preserves the same exact profile
\((a_k,b_k)\). A narrower replacement can change later residues
and continuation costs. The shared-table rule applies to the whole
profile and its realizing witness, not merely to a charge label.}

\subsection{L8. Sufficient finite bounds can be large}
\Plain{The proof supplies finite limits, but they can greatly exceed
the values in useful tables. A smaller limit needs a proof that it
cannot remove a required optimum. The incumbent construction's
charge is not a safe ceiling for every narrower-width query.}
\STE{The proof gives finite search limits. These limits can be much larger than the values in a useful table. A smaller limit can remove a table that gives the global optimum. Such a limit needs a separate proof.}
\Formal{The sufficient budget range is \(0\le C\le27k^3\), with
cutoff (15). At \(k=15\), the charge ceiling is \(91{,}125\).
The large \(S(k)\) is used below \(c_0\); the constructor bounds
thresholds at budgets at least \(c_0\).
Restricting all searches to \(C\le c_0\) can lose higher-charge,
lower-width alternatives. These bounds do not predict typical optima.}

\clearpage
\subsection{L9. Exact integers and memory constrain implementation}
\Plain{Row coefficients can exceed machine-integer ranges. Overflow,
approximate division or approximate equality can change feasibility.
Large state caches can exhaust memory. The theorem gives no practical
RAM estimate.}
\STE{The row coefficients require exact integer arithmetic. Fixed-width integers can overflow. An overflow can change a feasibility result. A state cache can require more memory than the computer has.}
\Formal{Coefficients under cap \(\Lambda\) have \(O(\Lambda)\) bits;
one matrix can require \(O(k^2(\Lambda+1))\) bits.
No fixed machine word suffices uniformly.
Division, valuation, equality and objective comparison must be exact.
Formula (9) removes the real-logarithm oracle.
A cache can contain many states, so its space is not bounded by
one state's encoding length.}

\subsection{L10. Faithful semantics are necessary}
\Plain{Small implementation changes can alter the problem. Negative
additions cost three. Both width histories must pass, including
peaks before division. Literal target signs and indexed identity
restoration must remain exact.}
\STE{The implementation must follow the specified operations and costs. It must do a check of both width histories before each normalization division. Each right shift must be exact. Each marker must match its target row. The final row matrix must be the identity matrix.}
\Formal{The transition relation is exactly (3)--(4), with raw
caps \(a\) and \(k-1+b\).
A normalized addition tests
\(1+\max(u_d+p_d,u_h+p_h)\) before subtracting its valuation.
Final \(I_k\) and exact-once literal target coverage are required.
Final-width checks alone are insufficient. Accepting
\(\pm v_k(w_\ell)\) interchangeably changes the domain.}

\clearpage
\subsection{L11. A partial search does not certify a minimum}
\Plain{A legal table establishes an achievable cost, not the absence
of a better answer. A timeout leaves optimality open unless another
complete certificate is available. Failure below a smaller cap does
not establish global infeasibility.}
\STE{A valid table gives an upper bound for the minimum cost. A timeout does not give a proof of optimality. A failed search below a smaller width limit does not show global infeasibility. The result must identify an incomplete search.}
\Formal{An incumbent of cost \(B\) proves only
\(\OPT(n,K)\le B\). Optimality also requires a valid lower bound
\(L\ge B\), supplied here by complete exact search.
A failed truncated query proves \(\theta(k,\delta,C)>n-2\),
not necessarily \(+\infty\). A hash or feasible-word replay does
not certify the absence of better words.}

\subsection{L12. The full implementation is not certified}
\Plain{The results are reviewed mathematical proofs with component
tests. The full optimizer is not machine formalized and was not
run at large arity in the reviewed package. The Fortran translation
needs its own validation and accurate completion-status reporting.}
\STE{The general proofs have independent mathematical reviews. The full algorithm has no complete machine-checked proof. The tests do not cover the full search domain. The Fortran implementation needs its own validation.}
\Formal{Executable correctness requires faithful implementation of
the transition relation, bounds, catalogue and exact outer solver.
Elementary formal lemmas and bounded checks do not prove the whole
program. They supply neither a full-domain runtime benchmark nor
a general compact optimality certificate.}

\clearpage
\section{Implementation result records}\label{sec:records}
This section specifies how to interpret results from the current
algorithm. It does not prescribe a new search strategy.
\begin{longtable}{@{}>{\raggedright\arraybackslash}p{0.26\textwidth}>{\raggedright\arraybackslash}p{0.68\textwidth}@{}}
\toprule Status & Required meaning\\ \midrule
Feasible incumbent&A replayed table or plan gives the stated achievable cost. No minimum claim follows.\\
Exact threshold&All necessary candidates for the query were covered; the threshold and witness are exact.\\
Infeasible within cap&No accepted word exists within the stated cap and budget. The scope must be recorded.\\
Globally infeasible threshold&The full sufficient saturation range was exhausted without a witness.\\
Certified recursive optimum&Complete local information and exact outer optimization establish (10)'s minimum for the declared \(n,K\) model.\\
Incomplete search&A time, memory, arithmetic-range or user cutoff prevented completion. A valid incumbent can still be reported.\\
\bottomrule
\end{longtable}

A reproducible record identifies \(n,K\), literal F.2 targets and
signs, charge weights, scanner rules, leaf formula and all cutoffs.
A local query also records \(k,\delta,C,\Lambda\), whether its cap
was truncated, its witness and completion status.

Practicality records should separate table-search time from outer
planner time. They should include visited prefixes or states,
maximum coefficient bit length, peak memory, exact-integer library,
compiler and build settings. These measurements describe an
implementation; they do not change its optimality domain.

\clearpage
\subsection{Controlled-English technical vocabulary}
The following technical nouns have one fixed meaning in the
controlled panels. Row, coefficient, matrix, integer, marker and
width history are defined in Section~\ref{sec:model}.
\begin{description}[style=nextline,leftmargin=0.35in]
\item[Arity] Number \(k\) of chunks in one recursive division.
\item[Charge] Weighted operation count \(\gamma\).
\item[Child width] Common integer width \(W_P(m)\).
\item[Complexity bound] Proved asymptotic bound on computation.
\item[Evaluation-point family] Fixed F.2 point sets indexed by arity.
\item[Feasibility] Satisfaction of the acceptance conditions.
\item[Global optimum] Minimum of (10) within its declared domain.
\item[Machine-checked proof] Proof accepted by a formal proof checker.
\item[Normalization division] Division by a row's maximal common power of two.
\item[Operand bitwidth] Number \(n\) or \(m\) of operand bits.
\item[Planner] Part that selects a recursive plan and shared tables.
\item[Recursive T-count] Cost function (10).
\item[Search domain] All permitted tables and recursive plans.
\item[State cache] Search states stored by an implementation.
\item[Table] Complete instruction sequence defined in Section~\ref{sec:dsl}.
\item[Upper bound] Value at least as large as the quantity being bounded.
\item[Validation] Checks that an implementation satisfies its specification.
\end{description}

\clearpage
\begin{thebibliography}{9}
\raggedright
\addcontentsline{toc}{section}{References and proof provenance}
\bibitem{kmy}
G. D. Kahanamoku-Meyer and N. Y. Yao.
\emph{Fast quantum integer multiplication with zero ancillas}.
arXiv:2403.18006v4, 14 November 2024.
Sections II--III; Appendix B, equations (24)--(27);
Appendices F.1--F.2 and G.
\url{https://arxiv.org/html/2403.18006v4}.
This source supplies notation, PhaseProduct and point-selection
context, not the local exhaustive-optimizer theorem.

\bibitem{ste}
ASD Simplified Technical English Maintenance Group.
\emph{ASD-STE100 Simplified Technical English}, Issue 9,
15 January 2025. Writing rules and dictionary.
\url{https://www.asd-ste100.org/STE_faq.html}.
Controlled panels use short sentences and consistent technical nouns;
independent conformity certification is not asserted.

\bibitem{kannan}
Ravi Kannan.
\emph{Minkowski's Convex Body Theorem and Integer Programming}.
Mathematics of Operations Research 12(3), 415--440, 1987.
Theorem 5.4, with the intermediate-integer bounds included in
the bit-complexity accounting.
\url{https://doi.org/10.1287/moor.12.3.415}.

\bibitem{threshold}
Local research package, 15 October 2026.
\emph{Constructive thresholds: exact FPT algorithm; uniform P remains open}.
These establish the finite local enumeration and cutoff.

\bibitem{outer}
Local research package, 15 October 2026.
\emph{Gap-closure result: polynomial outer optimization}.
Its earlier open local-realization obligation is supplied by
reference \cite{threshold}.

\bibitem{profile}
Local research package, 15 October 2026.
\emph{Exact three-integer planner profiles}.
Reference \cite{outer} replaces the earlier
outer-runtime bound of that note.

\bibitem{normalization}
Local research package, 15 October 2026.
\emph{Whole-table primitive normalization}.
The scanner-preserving replacement is incorporated into
reference \cite{threshold}.
\end{thebibliography}

\small
\textbf{Local sources.}
The local research packages cited above are internal working notes and are not
part of this repository.
This document does not claim uniform polynomial time for growing \(K\).
\end{document}

